\chapter{Data and Problem Analysis}
\label{chap:data}

\section{Dataset Overview}

This project studies drug--target interaction (DTI) prediction using benchmark
biomedical data. The main dataset used for exploratory analysis and baseline
experiments is \textit{Yamanishi08}. This dataset is suitable for the project
because it provides a gold-standard DTI network together with drug structure
features, protein sequence features, predefined cross-validation folds, and
knowledge graph triples.

The dataset can be viewed as a heterogeneous learning problem. On one side, the
DTI labels define a bipartite graph between drugs and target proteins. On the
other side, auxiliary feature sources describe the properties and biomedical
context of these entities. The goal of the machine learning model is to combine
these sources and predict whether a candidate drug--target pair represents an
interaction.

\begin{table}[H]
\centering
\caption{Main data sources in the Yamanishi08 dataset}
\label{tab:yamanishi_data_sources}
\begin{tabularx}{\textwidth}{lX}
\toprule
\textbf{File or folder} & \textbf{Description} \\
\midrule
\texttt{dt\_all\_08.txt} & Gold-standard positive DTI pairs. \\
\texttt{791drug\_struc.csv} & Drug identifiers and SMILES strings. \\
\texttt{morganfp.txt} & 1024-dimensional Morgan fingerprints for drugs. \\
\texttt{989proseq.csv} & Protein identifiers and amino acid sequences. \\
\texttt{pro\_ctd.txt} & 147-dimensional CTD descriptors for proteins. \\
\texttt{kg\_data/} & Biomedical knowledge graph triples. \\
\texttt{data\_folds/} & Predefined 10-fold train/test splits. \\
\bottomrule
\end{tabularx}
\end{table}

\section{Yamanishi08 DTI Network}

The file \texttt{dt\_all\_08.txt} defines the positive DTI network. Each row
contains a drug identifier, the relation type \texttt{DRUG\_TARGET}, and a
target protein identifier. Drug identifiers follow KEGG drug-style notation
such as \texttt{D00002}, while target identifiers follow KEGG gene-style
notation such as \texttt{hsa:10}.

\begin{table}[H]
\centering
\caption{Yamanishi08 DTI network statistics}
\label{tab:yamanishi_dti_stats}
\begin{tabular}{lr}
\toprule
\textbf{Statistic} & \textbf{Value} \\
\midrule
Known positive interactions & 5,127 \\
Unique drugs & 791 \\
Unique target proteins & 989 \\
Possible drug--target pairs & 782,299 \\
Interaction matrix density & 0.6554\% \\
Interaction matrix sparsity & 99.3446\% \\
Average targets per drug & 6.48 \\
Average drugs per target & 5.18 \\
\bottomrule
\end{tabular}
\end{table}

The interaction matrix is extremely sparse. Among 782,299 possible
drug--target pairs, only 5,127 are observed as positive interactions. This
means that most possible pairs are unobserved. In the supervised learning folds,
unobserved pairs are sampled as negative examples, but this should be interpreted
carefully: an unobserved pair is not always a biologically confirmed negative.
It may simply be an interaction that has not yet been reported.

\section{Degree Distribution and Network Sparsity}

The DTI graph is not uniformly connected. Most drugs and proteins have only a
small number of known interactions, while a small number of entities act as
hubs. In the Yamanishi08 network, 588 out of 791 drugs have at most five known
target proteins. This accounts for 74.34\% of all drugs. Similarly, 544 out of
989 target proteins interact with at most two drugs, accounting for 55.01\% of
all targets.

\begin{table}[H]
\centering
\caption{Low-degree entity statistics}
\label{tab:degree_sparsity}
\begin{tabular}{lrrr}
\toprule
\textbf{Entity type} & \textbf{Condition} & \textbf{Count} & \textbf{Percentage} \\
\midrule
Drugs & Degree $\leq 5$ & 588 / 791 & 74.34\% \\
Target proteins & Degree $\leq 2$ & 544 / 989 & 55.01\% \\
\bottomrule
\end{tabular}
\end{table}

This heavy-tailed structure has two important implications. First, models can
learn more reliable patterns for highly connected drugs and targets because
more interaction evidence is available. Second, prediction for low-degree
entities is more difficult because the model must rely more heavily on
descriptor features and graph context rather than direct interaction history.

\begin{table}[H]
\centering
\caption{Top hub entities in the Yamanishi08 DTI network}
\label{tab:dti_hubs}
\begin{tabular}{lrlr}
\toprule
\textbf{Drug} & \textbf{Degree} & \textbf{Target} & \textbf{Degree} \\
\midrule
\texttt{D02356} & 132 & \texttt{hsa:5743} & 61 \\
\texttt{D00528} & 113 & \texttt{hsa:5742} & 61 \\
\texttt{D01441} & 98 & \texttt{hsa:6331} & 36 \\
\texttt{D00294} & 83 & \texttt{hsa:154} & 35 \\
\texttt{D00437} & 81 & \texttt{hsa:1576} & 34 \\
\bottomrule
\end{tabular}
\end{table}

\section{Drug Structure Features}

Drug compounds are represented using two complementary sources. The file
\path{791drug_struc.csv} stores the SMILES representation of each drug, and
\path{morganfp.txt} stores the corresponding Morgan fingerprint vectors.
Morgan fingerprints are binary molecular descriptors that encode local
topological neighborhoods in the molecular graph.

\begin{table}[H]
\centering
\caption{Drug structure and fingerprint statistics}
\label{tab:drug_feature_stats}
\begin{tabular}{lr}
\toprule
\textbf{Statistic} & \textbf{Value} \\
\midrule
Number of compounds & 791 \\
Average SMILES length & 55.38 characters \\
Minimum SMILES length & 11 characters \\
Maximum SMILES length & 377 characters \\
Morgan fingerprint shape & $791 \times 1024$ \\
Average active bits per drug & 42.38 / 1024 \\
Fingerprint matrix sparsity & 95.86\% \\
\bottomrule
\end{tabular}
\end{table}

The fingerprint vectors are highly sparse: on average, only about 42 of the
1024 bits are active for each drug. This sparsity is useful because the features
are compact and interpretable as structural indicators. However, it also makes
generalization difficult for unseen drugs. In a drug cold-start setting, a test
drug may contain molecular substructures that are poorly represented in the
training set, causing descriptor-only models to struggle.

\section{Protein Sequence Features}

Target proteins are represented using amino acid sequences and CTD descriptors.
The file \texttt{989proseq.csv} stores protein sequences, while
\texttt{pro\_ctd.txt} stores 147-dimensional CTD descriptors. CTD features
summarize the composition, transition, and distribution of amino acid groups
under several physicochemical properties.

\begin{table}[H]
\centering
\caption{Protein sequence and CTD descriptor statistics}
\label{tab:protein_feature_stats}
\begin{tabular}{lr}
\toprule
\textbf{Statistic} & \textbf{Value} \\
\midrule
Number of target proteins & 989 \\
Average sequence length & 625.72 residues \\
Minimum sequence length & 83 residues \\
Maximum sequence length & 5,038 residues \\
CTD descriptor shape & $989 \times 147$ \\
CTD value range & 0 to 100 \\
\bottomrule
\end{tabular}
\end{table}

The CTD descriptor contains three groups of features. The composition and
transition components are typically in the range $[0,1]$, while distribution
features are expressed as positional percentages and can reach 100. This scale
difference motivates normalization before training models that are sensitive to
feature magnitude, especially neural networks and distance-based baselines.

\section{Knowledge Graph Data}

In addition to direct DTI labels and descriptor features, the dataset includes
biomedical knowledge graph triples. These triples connect drugs, proteins, genes,
pathways, functional annotations, domains, families, and other biological
entities. The graph provides relational context that can be used by KGE models
such as DistMult and CompGCN.

\begin{table}[H]
\centering
\caption{Knowledge graph statistics}
\label{tab:kg_stats}
\begin{tabularx}{\textwidth}{lrrrX}
\toprule
\textbf{KG file} & \textbf{Triples} & \textbf{Entities} &
\textbf{Relations} & \textbf{Main context} \\
\midrule
\texttt{kegg\_kg.txt} & 59,222 & 16,337 & 35 & KEGG pathways and genes \\
\texttt{yamanishi\_uniprot\_kg.txt} & 36,450 & 10,573 & 13 & UniProt annotations \\
\midrule
Combined KG & 95,672 & 25,487 & 48 & Biomedical graph context \\
\bottomrule
\end{tabularx}
\end{table}

The KEGG graph is dominated by pathway, gene, drug, compound, motif, and BRITE
relations. The UniProt-derived graph adds protein-centric context such as
biological process, cellular component, molecular function, protein pathway,
domain, family, and interaction relations. This motivates using knowledge graph
representation learning: the model can embed drugs and proteins according to
their wider biomedical neighborhoods, not only according to direct DTI labels.

\section{Cross-Validation Splits}

The dataset provides predefined 10-fold splits under four evaluation settings.
These settings are designed to test different levels of generalization.

\begin{table}[H]
\centering
\caption{Average fold composition across 10 folds}
\label{tab:fold_composition}
\small
\begin{tabular}{lrrrr}
\toprule
\textbf{Split} & \textbf{Train n} & \textbf{Test n} &
\textbf{Train +} & \textbf{Test +} \\
\midrule
Warm-start 1:10 & 44,864.2 & 5,538.9 & 4,615.0 & 512.0 \\
Warm-start 1:1 & 9,164.1 & 1,022.9 & 4,615.0 & 512.0 \\
Protein cold-start & 48,142.7 & 1,678.6 & 4,972.9 & 154.1 \\
Drug cold-start & 44,414.4 & 5,102.5 & 4,621.5 & 505.5 \\
\bottomrule
\end{tabular}
\end{table}

\begin{table}[H]
\centering
\caption{Average unique entity counts across evaluation splits}
\label{tab:split_overlap}
\small
\begin{tabular}{lrrrr}
\toprule
\textbf{Split} & \textbf{Train D} & \textbf{Test D} &
\textbf{Train T} & \textbf{Test T} \\
\midrule
Warm-start 1:10 & 791.0 & 223.9 & 989.0 & 985.2 \\
Warm-start 1:1 & 791.0 & 223.9 & 989.0 & 581.1 \\
Protein cold-start & 780.7 & 92.5 & 989.0 & 792.1 \\
Drug cold-start & 685.7 & 105.3 & 989.0 & 984.4 \\
\bottomrule
\end{tabular}
\end{table}

The warm-start settings evaluate held-out pairs while allowing drugs and
targets to appear in training. The 1:10 split is more imbalanced because it uses
approximately ten negative examples per positive example, while the 1:1 split is
approximately balanced.

For cold-start analysis, it is important to distinguish all sampled pairs from
positive interactions. In the drug cold-start split, positive test drugs do not
appear in positive training pairs, making it a strict test of generalization to
unseen drug structures. In the protein cold-start split, positive test targets
do not overlap with positive training targets. However, because negative samples
may contain many entities, counting all labeled pairs can still show entity
overlap. Therefore, cold-start conclusions should be based on the intended
positive-label partition rather than only raw entity overlap across all samples.

\section{Implications for Model Design}

The exploratory analysis suggests several important design requirements for the
proposed model:

\begin{itemize}
    \item \textbf{Handling sparsity:} The DTI matrix has only 0.6554\% observed
    positive density, so the model must learn from sparse supervision.

    \item \textbf{Using descriptor features:} Low-degree drugs and proteins
    have limited interaction history, making molecular fingerprints and CTD
    descriptors important for generalization.

    \item \textbf{Using graph context:} The knowledge graph contains 95,672
    triples and 48 relation types, providing relational evidence beyond direct
    DTI labels.

    \item \textbf{Supporting cold-start prediction:} Drug cold-start evaluation
    is especially challenging because unseen molecular scaffolds can produce
    out-of-distribution fingerprint patterns.

    \item \textbf{Modeling feature interactions:} The final pair representation
    combines sparse drug fingerprints, continuous protein descriptors, and dense
    KGE embeddings. This motivates the use of an NFM classifier, which is
    designed to capture interactions among heterogeneous features.
\end{itemize}

\section{Preprocessing Pipeline}

The preprocessing pipeline prepares the data for KGE training and supervised
DTI prediction:

\begin{enumerate}
    \item load the DTI fold files from \texttt{data\_folds/};
    \item load drug fingerprints and protein CTD descriptors;
    \item normalize descriptor features where required;
    \item load KG triples from KEGG and UniProt graph files;
    \item train the selected KGE model on graph triples;
    \item extract entity embeddings for drugs and target proteins;
    \item merge drug descriptors, protein descriptors, and KGE embeddings into
    pair-level feature vectors;
    \item train and evaluate supervised classifiers on each fold.
\end{enumerate}
